%% IJCARS submission.
%%
%% DOCUMENT CLASS: IJCARS uses Springer Nature's sn-jnl.cls, distributed from
%% springernature.com/gp/authors/campaigns/latex-author-support and not available on CTAN,
%% so it is not vendored in this repository. To switch to it, download the official
%% template (v3.1, Dec 2024), place sn-jnl.cls and the .bst files beside this file, and
%% replace the \documentclass line below with:
%%
%%     \documentclass[pdflatex,sn-mathphys-num]{sn-jnl}
%%
%% Everything else in this file is written to be compatible with that swap.
%%
%% Generated by src/build_latex.py -- do not edit by hand; edit the Markdown and rebuild.

\documentclass[11pt,a4paper]{article}

\usepackage[T1]{fontenc}
\usepackage[utf8]{inputenc}
\usepackage[margin=2.5cm]{geometry}
\usepackage{graphicx}
\usepackage{booktabs}
\usepackage{amsmath}
\usepackage{graphics}   % for \resizebox on wide tables
\usepackage[numbers,sort&compress]{natbib}
\usepackage[hidelinks]{hyperref}
%% microtype is omitted: not present in this TeX installation and purely cosmetic.

\graphicspath{{figures/}}
\setlength{\parskip}{0.4em}

\title{What does a surgical skill model measure? Quantifying identity confounding and
miscalibration on JIGSAWS}
\author{}
\date{}

\begin{document}
\maketitle

\begin{abstract}
\textbf{Purpose} JIGSAWS sanctions two cross-validation protocols for surgical skill assessment: one holds out a repetition while every surgeon remains in training, the other holds out a surgeon entirely. We ask what a model evaluated under the first measures, and whether reported confidence would support deferring uncertain trials to a human rater.

\textbf{Methods} Two prediction paths were evaluated on identical information across 103 trials, 8 surgeons and 3 tasks: 16 kinematic features with ridge and multinomial logistic regression, and a zero-shot commercial decision model given the same features rendered as text. Both protocols were run. Errors are reported against a feature-free reference predicting each held-out trial with its training-fold mean rating; intervals resample whole surgeons. Analysis rules were fixed in advance; post-hoc additions are labelled.

\textbf{Results} Changing only the protocol removed at least a third of skill-classification accuracy (paired gap 0.49-0.67, intervals excluding zero). A post-hoc diagnostic found the same features identified which of eight surgeons performed a held-out trial with accuracy 0.83-0.93 against chance near 0.13, separating identity confounding from underfitting. On unseen surgeons one task of three beat the reference, and there a single-feature regression on trial duration matched both models. Both paths were overconfident (confidence 0.755 against accuracy 0.379; 0.696 against 0.301), scoring worse than a constant forecast, and neither identified deferrable trials.

\textbf{Conclusion} Accuracy reported under the permissive protocol is consistent with surgeon recognition rather than skill measurement. With eight surgeons these are single-dataset demonstrations rather than established generalities. We recommend reporting held-out-surgeon results against an uninformative reference, with a subject-identifiability diagnostic.

\textbf{Keywords} Surgical skill assessment; JIGSAWS; identity confounding; cross-validation; calibration; selective prediction

\end{abstract}


\section{Introduction}

Automated assessment of surgical technical skill is motivated by the prospect of scalable, consistent feedback during training. If such a system is to inform training or credentialling decisions, it must generalise to surgeons who were not seen during its development. Whether current evaluation practice on the field's standard benchmark establishes that it does is the question we take up.

The concern is not specific to surgery. When a model is trained and tested on data from the same individuals, it can succeed by recognising the individual rather than the property of interest --- a failure mode named identity confounding \cite{chaibubneto2019detecting} and demonstrated systematically by Saeb et al. \cite{saeb2017need}, whose activity-recognition error stayed near 2\% when records from the same subjects appeared on both sides of the split, against roughly 7\% when whole subjects were held out and 27\% when only two subjects were available. Electroencephalography supplies the closest structural parallel to what we report below: an Alzheimer's classifier reaching 99.8\% accuracy under segment-based holdout fell to 53.0\% on unseen subjects, with a confidence interval spanning chance, and a seizure classifier fell from 79.1\% to 65.1\% \cite{brookshire2024leakage}. The same pattern appears in speech, where depression detection degrades sharply on unseen speakers because the acoustic features are entangled with speaker identity. These are instances of shortcut learning \cite{geirhos2020shortcut}, and leakage of this general kind has been documented across 17 fields, collectively affecting 294 papers \cite{kapoor2023leakage}. We return to the methodological debate this literature provoked in Section~2.3.

The benchmark itself ships two cross-validation protocols: one holds out a repetition while every surgeon remains in training, the other holds out a surgeon entirely. Both are sanctioned and both appear in this literature, so what follows is not a criticism of an informal practice. That performance falls under the second is not in dispute; it has been reported repeatedly. What has not settled is why. Three studies of JIGSAWS attribute the same gap to three different causes: limited generalisation, answered by a call for surgeon-invariant features; the quality of the dataset's representations and annotations; and an absence of subject-independent validation in earlier work. A fourth, on open-surgery hand tracking rather than robot kinematics, comes closest to the account we pursue, describing its own gap as quantifying ``the surgeon-specific component captured by video-level splits'' --- and stops there. Three different attributions, none of them measured, and a fourth study that names a mechanism without measuring it either. The test is inexpensive: train the same features to predict which surgeon performed the trial, and read off how much identity is recoverable. Such classifiers exist --- surgeons have been identified from intraoperative force profiles with high accuracy --- but they have been built to enable per-surgeon feedback, not to ask whether a skill benchmark is measuring skill. To our knowledge the test has not been run as a validity check on a skill-assessment protocol. The gap has been treated as a limitation to engineer away rather than as a quantity to measure.

Judging such a system also requires care about the target it is fitted to. The reference standard is a human rating, and human ratings of technical skill are neither perfectly reliable nor uniformly so: pooled inter-rater reliability for global rating scales in simulation-based assessment has been reported at 0.78 (95\% CI 0.71-0.83, 23 studies) \cite{ilgen2015checklists}, while untrained raters have been measured at 0.26, with the benefit of rater training lost during periods of non-use \cite{gawad2019interrater}. On this benchmark specifically, the self-declared skill labels track the expert ratings poorly \cite{hendricks2024exploring}. A model here is therefore fitted to a noisy and training-dependent judgment rather than to a fixed standard. Two consequences follow for how results should be read. Agreement statistics cannot be weighed against a single fixed bar. And a reported accuracy or correlation is not self-interpreting: it needs an explicit comparison with a predictor that uses no input at all, a point we return to in Section~3.5 with a concrete demonstration of how misleading its absence can be.

This paper asks what a skill model evaluated on JIGSAWS is actually measuring, and whether the confidence such a model reports would support the human-in-the-loop use these systems are proposed for. We evaluate two prediction paths on identical information --- a conventional baseline of 16 kinematic features with ridge and logistic regression, and a commercial zero-shot decision model given the same features rendered as text --- under both protocols the benchmark sanctions. We make four contributions. First, we quantify the protocol gap with paired surgeon-level intervals: changing only the evaluation protocol, with features, model and data held constant, removes at least a third of the reported skill-classification accuracy (Section~4.1). Second, in a post-hoc diagnostic, we show that the same features identify which of eight surgeons performed a held-out trial with accuracy between 0.83 and 0.93 against chance rates near 0.13, which distinguishes identity confounding from the underfitting explanation that a performance gap alone leaves open (Section~4.2). Third, we report what is to our knowledge the first calibration and selective-prediction audit of a surgical skill assessment model built from motion data, finding both paths overconfident on unseen surgeons and neither path's confidence able to identify the trials it gets wrong (Section~4.5, Section~4.6). Fourth, we benchmark both against a single-feature regression on trial duration, which matches them on the one task of three where any cross-surgeon signal survives (Section~4.4). We do not claim to have discovered identity confounding: the mechanism is well established elsewhere, and what we contribute is its demonstration, quantification and consequences in this setting.

\section{Related work}

\subsection{The benchmark and its two sanctioned protocols}

JIGSAWS \cite{gao2014jigsaws} is the standard public benchmark for robot-assisted surgical skill assessment, and the accompanying benchmark suite \cite{ahmidi2017dataset} defines the two cross-validation schemes used throughout this literature. Leave-one-supertrial-out holds out the \textit{i}-th repetition across all surgeons; leave-one-user-out holds out every trial from one surgeon. Both ship with the dataset and both are sanctioned, so the protocol we examine is the benchmark's own rather than an informal practice.

The field is aware that individual style exists. A review of gesture recognition in robotic surgery \cite{vanamsterdam2021gesture} notes that leave-one-user-out is the harder setting ``due to subject-specific style variability'', and surveys methods that normalise such variability away. That framing treats identity as nuisance variance to be suppressed in pursuit of a better model. Our concern is the reverse: the same variability, under a protocol that places a surgeon on both sides of the split, does not degrade a score but inflates it.

\subsection{Label validity and protocol validity are separate critiques}

The closest existing critique of skill assessment on this benchmark concerns the target. Hendricks et al. \cite{hendricks2024exploring} show that the self-declared novice/intermediate/expert label correlates poorly with the expert GRS rating, and \cite{hoffmann2026looking} reports the same mismatch. Our own data reproduce it --- one surgeon labelled expert averages GRS 13.5 while one labelled novice averages 19.1 (Section~3.1) --- and we report that as a replication of their finding rather than as one of ours.

The distinction from our contribution is worth stating precisely, because the two are easily conflated. \textbf{Hendricks et al. ask whether the target measures skill. We ask whether the evaluation protocol permits a model to succeed without measuring skill at all, whatever the target's quality.} These critiques are independent, and they compound. A perfectly valid target would not repair what we describe: surgeon identifiability is a property of the feature representation and the fold structure, not of what the label means, so a model could shortcut a flawless label just as readily. Conversely their critique survives a flawless protocol, since a well-designed split cannot rescue a target that measures the wrong thing. A reader persuaded by Hendricks et al. should not conclude that fixing the label would fix the benchmark.

\subsection{The performance gap is widely reported, and framed three other ways}

That accuracy falls under leave-one-user-out is not news in this literature. What differs is the explanation offered.

An interpretable whole-stream video framework \cite{chen2026interpretable} reports near-perfect supertrial-out accuracy across the three tasks, a significant drop under user-out evaluation, and near-chance transfer to an external dataset, concluding that generalisation to unseen surgeons is limited and calling for ``surgeon-invariant feature learning''; the goal is named, the mechanism is not measured. \cite{hoffmann2026looking} evaluates GRS prediction under a strict disjoint-participant protocol and reports a concordance correlation of -0.03 (+/-0.08) for a supervised model trained and evaluated entirely on JIGSAWS, below chance, against 0.22 for its best baselines; having ruled out model-side causes through an extensive architecture and hyperparameter search, it attributes the failure to the annotations, which it reports ``lack a consistent ordering with respect to the visual features that encode actual skill differences''. A recent explainable-machine-learning study using rich kinematic features \cite{xuan2026explainable} criticises prior JIGSAWS work for lacking subject-independent validation --- the same complaint, again without a direct test. Lefor et al. \cite{lefor2020motion} earlier flagged the dataset's subject imbalance and an inverse correlation among subjects as a threat to generalisation, without quantifying either.

Closest to our own framing is the markerless hand-tracking study of open-surgery suturing by Constans et al. \cite{constans2026markerless}, which finds accuracy falling to 52.3\% under leave-one-subject-out and states that this gap ``quantifies the surgeon-specific component captured by video-level splits''. We should be plain about how little separates that sentence from our thesis: they articulate the same interpretation of the same kind of gap, and had they run a subject-identification probe, the attribution would already exist in this literature. Our contribution relative to that study is correspondingly narrower than relative to the others above --- a different modality, robot kinematics rather than open-surgery hand tracking, and a direct measurement of the component they name but do not test. Taken together these studies establish that the field suspects a subject-specific confound; none tests it directly.

The null in ``Looking Beyond the Scale'' overlaps our Section~4.3, and that overlap strengthens rather than duplicates the contribution. Their study and ours reach the same conclusion about cross-participant GRS prediction through almost entirely different routes: different feature representations, different model families, a different agreement metric, and a different reference against which the result is judged. A single null is always vulnerable to the objection that it reflects an idiosyncrasy of one pipeline --- an underpowered search, an unlucky architecture, a poorly chosen metric. Two independent pipelines arriving at the same place is much harder to explain that way, and their architecture search excludes model-side causes that our single linear model could not have excluded alone. Independent corroboration by a different method is better evidence than either result on its own. What their study does not supply, and what we add, is the attribution: they establish that the signal is absent, while we measure where the apparent signal under the permissive protocol comes from.

One counter-argument must be met directly, because it applies to every study above and to our Section~4.1. Responding to Saeb et al. \cite{saeb2017need}, Little et al. \cite{little2017using} argued that a gap between record-wise and subject-wise cross-validation can arise from \textit{underfitting} on heterogeneous data rather than from identity confounding, so that a gap alone establishes nothing about mechanism. The objection is sound, and none of the surgical studies above can answer it, because each reports only a gap. Our Section~4.2 is designed to answer it: the same 16 features, the same model and the same folds identify which of eight surgeons performed a held-out repetition with accuracy between 0.83 and 0.93, against chance rates near 0.13. A performance gap is ambiguous between confounding and underfitting; recoverable identity at roughly seven times chance is not.

Surgeon identifiability itself is not a new observation. Working from intraoperative tool-tissue force data rather than benchmark kinematics, Baghdadi et al. \cite{baghdadi2023force} identify individual surgeons with micro-average AUC 0.97 and 83.8\% accuracy across 89 neurosurgical cases, concluding that each surgeon has a unique force profile recoverable by deep networks. They frame that finding as a capability: a route to bespoke per-surgeon feedback and to correlating outcomes with individual performance. It is not connected there to the validity of any skill-assessment protocol, and no performance gap is reported.

Our claim is correspondingly narrow. We are not the first to identify surgeons from surgical signals; what we add is the use of that identifiability as a validity diagnostic, run on the same features, model and folds as the skill task, and tied to the protocol gap it explains. To our knowledge no study both attributes that gap to surgeon-identity leakage and demonstrates the leakage by classifying the surgeon from the features used for skill. This is the outcome of a recorded search of PubMed and general web sources (\texttt{paper/systematic\_search.md}); Google Scholar and IEEE Xplore could not be queried programmatically, so the claim is bounded rather than exhaustive.

\subsection{Reported performance, and the missing reference}

In the studies we examined \cite{zia2018automated,wang2021towards,vbanet2022video}, results are reported as accuracies or as correlations between predicted and expert scores, without the value the same metric would take for a predictor using no motion information at all. We have not surveyed this literature systematically enough to call that a field-wide norm; it is the pattern in the work we read. We depart from it deliberately, for the reason given in Section~3.5 and Section~3.7: under leave-one-user-out, a predictor that returns its training fold's mean GRS scores approximately -0.8 on pooled rank correlation purely by construction. A positive correlation is therefore not self-interpreting, and a model can post one while committing larger absolute errors than that uninformative predictor.

\subsection{Calibration and selective prediction: a borrowed lineage}

Calibration has begun to receive attention in surgery: \cite{surgiq2026} benchmarks large language models on surgical knowledge questions and includes a calibration analysis, reporting that even strong models make confident mistakes on plausible distractors. That is a different setting from ours. We located no calibration or selective-prediction study of a surgical \textit{skill assessment} model built from motion, kinematic or video data, and surgical work on confidence-gated deferral is correspondingly thin. The methodological lineage we draw on is therefore from diagnostic imaging and intensive care, where deferral is established and demonstrably useful: CoDoC \cite{dvijotham2023codoc} reduces false positives by 25\% at matched sensitivity while cutting clinician workload, and ``Trust or Escalate'' \cite{zheng2024trust} provides calibrated selective evaluation with provable agreement guarantees. Overconfidence in model-as-judge settings is itself an active topic \cite{tian2025overconfidence}. We take from this literature one caution that bears directly on our own null result: deferral cannot manufacture reliability that the underlying model lacks \cite{calibratedtriage2026}.

\subsection{Zero-shot and text-based scoring}

Foundation models have reached this benchmark, but not in the role we examine. The nearest work, ZEAL \cite{zeal}, uses zero-shot video segmentation to produce features for a \textit{trained} LSTM, so the learning still happens downstream; that is a different mechanism from a decision model scoring a rubric with no training examples at all. The model we evaluate, Jev, is marketed on calibration --- its training method is branded Reinforcement Learning for Calibrated Decisions --- though a trade critique notes the term is coined and that its ``zero hallucinations'' claim concerns schema conformance rather than correctness \cite{eliot2026forbes}. Rather than rely on vendor claims, we note three independent evaluations published in September 2026 \cite{priorbench2026,scienthoon2026,webofmike2026} converging on a consistent pattern: calibration holds on in-distribution data and degrades markedly out of distribution, and one reports that the model's separate confidence field is less informative than its maximum class probability. We give no figures from these evaluations, for the reason that follows. \textbf{All three are non-peer-reviewed single-author repositories or blog posts, and we cite them as preliminary external evidence rather than settled results}; they also identify the model by version strings that may not correspond exactly to the build we pinned, so we treat them as evidence about the model family rather than as replications of our setup. The pattern they describe is the one our Section~4.5 reports on unseen surgeons.

\section{Methods}

\subsection{Dataset}

See Table 1.

We use JIGSAWS, a public dataset of robot-assisted surgical training trials recorded on a da Vinci system at 30 Hz. Eight surgeons performed three bench-top tasks --- suturing, knot-tying and needle-passing --- with up to five repetitions each. Each trial carries an expert rating on a modified OSATS instrument: six elements scored 1-5, whose sum is the global rating score (GRS), range 6-30. Each surgeon also carries a self-declared skill label (novice, intermediate, expert) reflecting hours of robotic experience.

The per-task metadata files, rather than the directory listing, are authoritative for which trials are usable: the kinematics directories contain 40 files per task, but the metadata files list 39 suturing, 36 knot-tying and 28 needle-passing trials, 103 in total. The surplus recordings are excluded, and one of them is empty. Surgeon G contributes no usable needle-passing trial, so that task has seven surgeons rather than eight. We verified that GRS equals the sum of its six elements on all 103 rows, with no exceptions, which is what identifies the columns; the ordering of the six elements among themselves rests on the dataset's own readme and cannot be confirmed arithmetically, so we report no per-element result.

Per-task GRS standard deviations are 5.40 (suturing), 5.11 (knot-tying) and 4.82 (needle-passing). These set the scale against which prediction errors below should be read. We note that the self-declared skill label tracks the expert GRS rating poorly: one surgeon labelled expert has a mean GRS of 13.5 while one labelled novice has 19.1 (\texttt{addendum\_label\_vs\_grs}). The two targets are therefore reported separately throughout rather than treated as interchangeable.

% PROVENANCE: Source: dataset facts are enforced by assertions in src/loaders.py and its tests rather
than logged as experiments.*

\subsection{Features}

We compute 16 descriptive motion features per trial from the patient-side manipulator channels --- the instruments inside the workspace --- rather than the surgeon-side console channels. Features are computed over the annotated task window, defined as the first gesture start to the last gesture end from the trial's transcription; this window matches the one embedded in the official split filenames on all 39 suturing trials.

Tool-tip positions are low-pass filtered once, before any derived quantity, with a zero-phase Butterworth filter at 5 Hz (second order per pass, applied forward and backward, so the gain at the cutoff is 0.5 and no time shift is introduced). The cutoff was selected to sit below the 8-12 Hz physiological tremor band and above the low-frequency content of voluntary movement; no JIGSAWS-specific value exists in the literature, so the choice is a judgment within a defensible range and is accompanied by sensitivity runs at no filtering and at 8 Hz.

The feature set comprises total task time; and per instrument, path length, mean and standard deviation of speed, mean acceleration magnitude, the fraction of time below a 1 mm/s motion threshold, convex-hull working volume, and the number of gripper open-close cycles; plus the correlation between the two instruments' speed profiles. Speed is computed by finite differences of position rather than from the recorded velocity channels, so it can be validated against hand-computed synthetic trajectories; on one trial the two agree at Pearson r = 0.88 and 0.93 for the two instruments, with finite-difference means within 5\% of the recorded means (\texttt{addendum\_speed\_validation}).

All thresholds were fixed in \texttt{configs/features.yaml} before use and were never adjusted against any outcome. Three features that appeared in our initial plan were removed before modelling, each for a stated reason unrelated to predictive performance: maximum speed and maximum acceleration, because both are corrupted by single-frame position discontinuities (the largest is 26.9 mm within one 1/30 s frame against a typical 0.2 mm, and 47 of 206 arm-trials contain a step above 5 mm; \texttt{addendum\_position\_jumps}); spectral arc length, because whole-trial values correlate -0.83 with movement count and -0.81 with trial duration and therefore do not isolate smoothness (\texttt{addendum\_sparc\_confound}); and movement count, because threshold crossings remain unstable after filtering.

% PROVENANCE: Source: configs/features.yaml; feature-matrix construction is tested in
\texttt{tests/test\_features.py}.*

\subsection{Two prediction paths}

\textbf{Baseline.} Ridge regression predicts GRS and multinomial logistic regression predicts the three-way skill label. Features are standardised inside a pipeline so that scaling parameters are estimated from training trials only. The ridge penalty is selected within each training set by leave-one-subject-out search over \{0.1, 1, 10, 100\}, which prevents the inner search from rewarding a model that recognises an individual surgeon's style. The logistic regulariser is fixed at C = 1.0 rather than searched: with two experts and two intermediates, a subject-level inner search can leave a training set containing no instance of a class, making the search ill-defined. This asymmetry between the two models follows from the data rather than from preference.

\textbf{Zero-shot decision model.} The second path uses Jev (TypeSafe, model version pinned to \texttt{jev-1.13.0}), a commercial model that answers typed questions against a supplied state and returns a score with a probability per level and a confidence value. It is never trained. Because the model's documentation states that it does not count reliably and cannot judge whether two numeric values are close, no number is transmitted: each of the 16 features is rendered in code as one of five ordinal phrases describing its percentile within the same task (``typical of comparable attempts'', ``higher than most comparable attempts'', and so on). The state contains no trial identifier, no surgeon letter, no repetition index, no label and no digit; tests enforce each of these conditions. Seven scoring questions --- one overall and six OSATS elements, each with five ordered level descriptions --- are sent in a single request per trial. Each trial was scored three times, giving 309 requests. Every raw response is cached to disk before parsing, and the model version reported in each response is checked against the pinned version on every read.

The rubric wording was authored before any trial was scored and was not revised afterwards. It was, however, written after we had seen the feature distributions and the baseline results, and it refers to completion time and pauses; this is a disclosed route by which the model's output could come to resemble a duration measure, and it bears on the interpretation in Section~4.4.

% PROVENANCE: Source: configs/baseline.yaml, configs/jev.yaml; kind=regression,
\texttt.*

\subsection{Evaluation protocol}

JIGSAWS defines two cross-validation schemes. Leave-one-supertrial-out (LOSO) holds out the \textit{i}-th repetition across all surgeons, so every surgeon appears in both training and test data. Leave-one-user-out (LOUO) holds out all trials from one surgeon, so the test surgeon is unseen. Pre-computed skill-detection splits ship only for suturing; for the other two tasks we generate folds by the same documented rule. That generator reproduces the official suturing folds exactly --- identical training and test membership in all five LOSO and all eight LOUO folds --- which is verified by a test, so the generated folds instantiate the standard definition rather than approximating it.

Every trial is predicted exactly once per scheme, by a model fitted without it. We report both schemes throughout. The contrast between them is the paper's primary object of study.

% PROVENANCE: Source: src/loaders.py; fold equivalence is pinned by
\texttt{test\_generated\_splits\_match\_official\_suturing}.*

\subsection{Reference predictor and the headline error metric}

Predictive claims are stated against an explicit reference rather than against zero. The \textbf{feature-free reference} predicts every held-out trial with the mean GRS of its own training fold; it uses no motion information whatsoever. A model has demonstrated something only if its error falls below this reference's error.

Our headline quantity is therefore the \textbf{reduction in mean absolute error relative to that reference}, in GRS points, positive meaning the model is better. This metric replaced the one we pre-registered, for a reason given in Section~3.7.

\textbf{Interpretation rule, fixed before any interval was computed:} a method beats the reference only if the lower bound of the 95\% interval on its MAE reduction lies above zero. With eight surgeons this is a weak test, and we treat ``excludes zero'' as suggestive rather than as proof; no analysis selects models by it.

To place Jev's ordinal score on the GRS scale, a single ordinary-least-squares line is fitted within each training fold and applied to the held-out trials, so the calibration uses no held-out information. Because a fitted line would silently rescue a score that ranked skill backwards, we record the sign of the fitted slope in every fold as a direction guard.

% PROVENANCE: Source: kind=headline_mae_vs_reference, kind=jev_headline_mae_vs_reference,
\texttt{kind=diagnostic\_feature\_free\_reference}.*

\subsection{Uncertainty}

Trials from one surgeon are not independent, so all intervals come from a bootstrap that resamples \textbf{whole surgeons} with replacement (2000 resamples, seed 0), never individual trials. Predictions are held fixed during resampling. An interval therefore expresses which surgeons happened to be evaluated; it does \textbf{not} express how a result would move if the model were refitted on different training data, and it does not capture that the eight surgeons are themselves a small and possibly unrepresentative sample. Resamples in which a statistic is undefined are counted and reported rather than silently dropped; for every regression result reported here that count is zero. For the paper's central claim we additionally bootstrap the LOSO-minus-LOUO difference \textbf{paired on the same resampled surgeons}, which is more informative than comparing two independently computed intervals.

% PROVENANCE: Source: src/metrics.py; the whole-surgeon property is pinned by
\texttt{test\_bootstrap\_resamples\_whole\_surgeons\_never\_individual\_trials}.*

\subsection{Protocol scrutiny}

Two findings about our own protocol belong in the record, and they are distinct.

\textbf{Leakage was tested for and was not found.} Scaling parameters and hyperparameters are estimated inside training folds only, and held-out trials are never seen by any fitted object. Two tests verify this behaviourally rather than by inspection: replacing the labels of a test fold with nonsense leaves that fold's predictions bit-identical, and altering some test-fold features leaves the predictions for the remaining test trials unchanged. To confirm the tests can detect the failure they are designed for, a deliberately leaky variant --- standardising all rows before splitting --- moved untouched test-fold predictions by 16 GRS points where the correct pipeline moved them by exactly zero. We additionally refitted the regression 50 times per task and scheme under shuffled labels; every mean rank correlation was negative, which is the opposite of the direction leakage produces.

\textbf{Separately, our pre-registered headline statistic proved uninterpretable under LOUO and was replaced.} We had pre-registered the pooled out-of-fold Spearman correlation. Under LOUO each fold contains one surgeon and the prediction is one constant per fold; surgeons differ in mean GRS, and a surgeon with a high mean receives a lower training-fold mean, so prediction and truth move in opposite directions across folds by construction. Measured, the feature-free reference scores -0.78 to -0.86 on this statistic. A LOUO correlation therefore cannot be read against zero. We report the MAE reduction of Section~3.5 as the headline and retain the correlations, always beside the reference value, never alone.

No result was altered by either finding. The first is a negative check that passed; the second is a change of metric, and it is a post-hoc change (Section~4, preamble). We also note that the shuffled-label criterion we wrote was two-sided and so flagged two of six cells in the harmless negative direction; it was replaced with a one-sided criterion in a second configuration, with the original flags left on the record.

% PROVENANCE: Source: kind=permutation_sentinel, kind=diagnostic_intercept_only; leakage behaviour
is pinned by tests in \texttt{tests/test\_baseline.py}.*

\subsection{Pre-registration}

Analysis rules for each phase --- feature thresholds, both evaluation schemes, model families, the metrics, the bin counts, the coverage levels, and the interpretation rule of Section~3.5 --- were committed to a version-controlled log before the corresponding analysis was run. Where an analysis was instead added after seeing results, it is labelled as such in Section~4. Two such additions are material and are identified there: the surgeon-identifiability diagnostic, and the metric substitution of Section~3.7.

---

\section{Results}

Every subsection below is labelled \textbf{[pre-declared]} where the analysis was specified before the results were seen, or \textbf{[post-hoc]} where it was added afterwards to explain something already observed. In summary: the protocol comparison (Section~4.1), the calibration analysis (Section~4.5) and the selective-prediction analysis (Section~4.6) were pre-declared; the surgeon-identifiability diagnostic (Section~4.2), the task-time comparison (Section~4.4) and the headline error metric itself (Section~4.3) were not. The distinction is load-bearing for a paper that examines protocol discipline, so it is carried in the text and not only in a table.

\subsection{The protocol gap [pre-declared]}

Skill classification accuracy under LOSO, where every surgeon appears in training, is 0.92 (suturing, 95\% interval 0.85-0.98), 0.97 (knot-tying, 0.92-1.00) and 0.93 (needle-passing, 0.85-1.00). Under LOUO, where the test surgeon is unseen, the same features and the same model give 0.44 (0.25-0.63), 0.31 (0.08-0.58) and 0.39 (0.16-0.69).

The drop is the paper's central observation, and it is best expressed as a paired difference. Bootstrapping the LOSO-minus-LOUO gap on the same resampled surgeons gives +0.49 (0.35-0.62) for suturing, +0.67 (0.41-0.89) for knot-tying and +0.54 (0.27-0.71) for needle-passing. All three intervals exclude zero by a wide margin. Changing only the evaluation protocol, with the features, the model and the data held constant, removes at least a third of the reported accuracy, and plausibly more than half; the point estimates are 0.49 to 0.67 but the intervals reach down to 0.27.

The LOUO point estimates fall at or below the rate obtainable by always predicting the most frequent class in that task --- 0.487, 0.444 and 0.393 respectively, so suturing and knot-tying fall below and needle-passing lands exactly on it. \textbf{We state this as a point estimate and not as a demonstration.} The LOUO intervals are wide enough to span the majority-class rate in all three tasks, so with eight surgeons we cannot establish that new-surgeon accuracy is below that floor; we can establish only that it is far below the LOSO figure. Quadratic-weighted kappa is negative under LOUO for suturing (-0.19) and needle-passing (-0.47), indicating agreement worse than chance on those two tasks, against +0.77 and +0.80 under LOSO.

% PROVENANCE: Figure 1A; Table 2. Source: kind=addendum_classification_ci, kind=addendum_protocol_gap,
\texttt{kind=addendum\_majority\_class}, \texttt (Table T2).*

\subsection{The same features identify the surgeon [post-hoc]}

This analysis was added after observing the gap in Section~4.1, in order to explain it. It was not part of the pre-registered plan, and it should be read as a diagnostic rather than as a confirmatory test.

Substituting surgeon identity for skill as the target, and leaving the features, the model, the folds and the regularisation unchanged, the same 16 features identify which of eight surgeons performed a held-out repetition with accuracy 0.87 (95\% interval 0.77-0.97) for suturing, 0.83 (0.67-0.97) for knot-tying and 0.93 (0.83-1.00) for needle-passing, against chance rates of 0.125, 0.125 and 0.143. All three intervals lie far above chance.

Because the skill label is a property of the surgeon and is constant across that surgeon's trials, a model that can identify the surgeon can retrieve the label without representing skill. This is a sufficient mechanism for the LOSO figures in Section~4.1, and it is consistent with their collapse under LOUO, where identity carries no usable information about an unseen individual. It is not proof that the classifier relies on identity: we have demonstrated that the information is present and recoverable, not that this particular model uses it. Distinguishing the two would require an intervention on the features, which we did not perform.

% PROVENANCE: Figure 1B; Table 2. Source: kind=addendum_surgeon_identity_ci,
\texttt{kind=diagnostic\_surgeon\_identity}.*

\subsection{GRS prediction against a feature-free reference [pre-declared design, post-hoc metric]}

The design of this analysis was pre-registered; the metric reported was not, for the reason given in Section~3.7. We flag that explicitly because the substitution is favourable to readability and could otherwise pass unnoticed.

Under LOUO, ridge regression reduces mean absolute error relative to the feature-free reference by +1.54 GRS points (95\% interval 0.65-2.52) on knot-tying, from 5.10 to 3.56. On suturing the reduction is +0.15 (-1.04 to 1.34) and on needle-passing it is -0.50 (-1.21 to -0.04) --- that is, on needle-passing the fitted model is \textit{worse} than predicting the training mean, with an interval excluding zero on the wrong side. Under the pre-declared rule, only knot-tying beats the reference. The zero-shot path behaves the same way: +2.30 (1.14-3.45) on knot-tying, -0.05 (-1.05 to 0.96) on suturing and -0.24 (-0.79 to 0.31) on needle-passing. The two knot-tying reductions are not distinguishable from one another: their intervals overlap across almost their whole length, and nothing in this study ranks the two paths. Its fitted slope was positive in all 38 folds, so the direction guard of Section~3.5 passed and no inverted score is being rescued by the calibration line.

Under LOSO the same comparison is positive on all three tasks for ridge (+1.15, +1.78, +0.43, all intervals excluding zero), which is the Section~4.1 pattern reappearing in the regression target.

Agreement statistics tell a consistent story. Under LOUO, ICC(2,1) is 0.60 (0.35-0.78) for ridge and 0.73 (0.53-0.85) for the zero-shot path on knot-tying, and is indistinguishable from zero on the other two tasks for both paths (ridge 0.20 and 0.00; zero-shot 0.06 and -0.13). Bland-Altman bias is small everywhere (-0.37 to +0.52 GRS points), but the limits of agreement span roughly ±11 points on suturing and needle-passing against ±7 on knot-tying --- that is, on two of three tasks the disagreement between model and rater is about as large as the full spread of the scale. For context, pooled inter-rater reliability for global rating scales in simulation-based assessment has been reported at 0.78 (95\% CI 0.71-0.83, 23 studies; Ilgen et al. 2015), with the original OSATS report giving 0.64-0.72 \cite{martin1997osats}. Neither source names the coefficient it pooled, so neither is strictly an ICC(2,1) and both comparisons are approximate. Our strongest agreement, 0.73 for the zero-shot path on knot-tying, sits below the pooled figure and carries an interval (0.53-0.85) wide enough to be uninformative about the comparison. We avoid the language of a ceiling: rater reliability is not a fixed bar but depends heavily on training, with untrained raters reported at ICC 0.26, rising to 0.76 immediately after training but varying widely on later re-testing (-0.17 to 0.70), which the authors summarise as the benefit of training being ``lost in periods of nonuse'' \cite{gawad2019interrater}. The target these models are fitted to is therefore itself noisy, and the reading we take is simply that no method here demonstrated agreement comparable to trained human raters on any task.

We note that the pooled out-of-fold Spearman correlations for these same predictions (LOUO: 0.21, 0.64, -0.13 for ridge) must be read against the feature-free reference's values on the identical folds (-0.78, -0.86, -0.79) rather than against zero, for the reason established in Section~3.7.

% PROVENANCE: Figure 1C; Table 3. Source: kind=headline_mae_vs_reference, kind=jev_headline_mae_vs_reference,
\texttt{kind=eval\_agreement}, \texttt *

\subsection{A one-feature comparison [post-hoc]}

This comparison was added after Section~4.3, to ask what the surviving knot-tying signal consists of. It is diagnostic.

The zero-shot model's overall score correlates with total task time at Spearman -0.91 (suturing), -0.81 (knot-tying) and -0.86 (needle-passing). Correlations with all 16 features are logged, so this is not a selected maximum: task time is the strongest associate in every task, with the next strongest between 0.41 and 0.71 in absolute value.

Substituting a single ordinary-least-squares line on task time alone for either model, on identical folds and under the identical metric, reduces error by +1.88 GRS points (0.94-2.89) on knot-tying under LOUO. On suturing (+0.10, -0.72 to 0.83) and needle-passing (-0.29, -0.88 to 0.11) it fails to beat the reference, as both models do. The three knot-tying intervals --- 0.65-2.52 for ridge, 1.14-3.45 for the zero-shot path and 0.94-2.89 for task time alone --- overlap across almost their whole length. We therefore report them as three estimates of one effect and draw no ordering among them; the point estimates differ, but this study cannot resolve that difference.

On this evidence the surviving cross-surgeon signal is, to within the resolution of eight surgeons, the signal available from trial duration alone. We do not claim the two are equivalent; we claim that this study does not demonstrate any advantage for either the 16-feature model or the zero-shot model over a one-feature baseline on the one task where a signal survives.

% PROVENANCE: Source: kind=addendum_jev_score_vs_feature, kind=posthoc_time_only_headline.
\subsection{Calibration [pre-declared]}

Each path is assessed against the target it was given: the baseline classifier against the self-declared skill label it was trained on, and the zero-shot model against the level implied by GRS, raw and with no recalibration applied. The two are therefore not compared with one another.

On a held-out surgeon, the baseline classifier states a mean confidence of 0.755 while achieving 0.379 accuracy. Its expected calibration error is 0.377 (0.237-0.541), and its Brier score of 0.875 is worse than the 0.647 obtained by a constant forecast equal to the observed class frequencies --- a Brier skill score of -0.351, with an interval (-0.688 to -0.092) lying entirely below zero. Its most confident bin, containing 46 of 103 trials at a mean stated confidence of 0.92, is correct 39\% of the time (\texttt{addendum\_reliability\_bin}).

Under LOSO the same classifier appears well calibrated and, if anything, underconfident: mean confidence 0.837 against accuracy 0.942, expected calibration error 0.105, Brier skill +0.778 (0.698-0.856). The same model and the same features thus yield opposite calibration verdicts under the two protocols, which is Section~4.1 restated in the confidence domain.

The zero-shot model is also overconfident: mean confidence 0.696 against accuracy 0.301, expected calibration error 0.395 (0.317-0.494), Brier skill -0.433 (-0.588 to -0.242). Its predictions are usually close --- within one level on 84\% of trials --- but they sit systematically high: the mean predicted level is 2.05 against a mean true level of 1.71, an offset of about a third of a level, and it places 39 trials in the top two levels where the GRS-implied truth places 22 (\texttt{addendum\_jev\_level\_agreement}, \texttt{addendum\_jev\_level\_usage}). That offset means its exact-level miscalibration mixes genuine overconfidence with a scale mismatch between its rubric and our GRS-derived levels; the two cannot be separated without a recalibration that was excluded in advance, and we therefore do not attribute the full -0.433 to overconfidence.

% PROVENANCE: Figure 2 (eval_reliability.png); Table 4. Source: kind=eval_calibration,
\texttt{kind=addendum\_mean\_confidence}, \texttt{kind=addendum\_reliability\_bin}.*

\subsection{Selective prediction [pre-declared]}

The deployment question is whether a model's confidence identifies the trials on which it should defer to a human rater. We rank the 103 pooled trials by each path's own confidence, defer the least confident first, and measure the absolute GRS error of those retained.

\textbf{The two paths do not support the same inference here, and we state that before the numbers.} Ridge regression emits no native confidence, so some proxy was required; we declared in advance that the baseline would pair its regression prediction with the classifier's top-class probability. Those are two different models fitted to two different targets, and Section~3.1 establishes that the targets disagree --- one surgeon labelled expert averages GRS 13.5 while one labelled novice averages 19.1. The baseline's confidence therefore carries no information about the regression's error except whatever the two targets happen to share. \textbf{Its null result is uninformative about whether ridge knows when it is wrong; it is a property of the pairing.} The zero-shot path has no such problem: its score and its confidence come from one model answering one question, so it is the only clean test of self-knowledge in this study.

The clean test is null. Ordering by the zero-shot model's own confidence lowers the area under the risk-coverage curve by 0.155 GRS points relative to random ordering under LOUO, with a 95\% interval of -0.416 to 0.606; substituting its maximum level probability for its reported confidence changes nothing (0.149). Under LOSO the figure is 0.077. The uninterpretable arm is also null, for the record: 0.135 (-0.628 to 0.922) under LOUO and 0.045 under LOSO. Expressed as a fraction of the distance between random and ideal ordering these are 7.4\% and 5.5\% under LOUO, and 3.9\% and 2.5\% under LOSO --- percentages of a quantity that is itself indistinguishable from zero, so they should be read only alongside the intervals above and not quoted alone.

The coverage table shows the same flatness directly. Deferring 20\% of trials by baseline confidence under LOUO changes the retained error from 4.31 to 4.50 GRS points --- it rises --- and deferring 60\% gives 4.15. The zero-shot path falls from 4.05 to 3.76 at 40\% coverage, a change no interval supports.

One further caution. The upper interval bounds are +0.61 GRS points for the clean arm and +0.92 for the other, up to a fifth of the overall error, so this is an under-powered null rather than a demonstration that confidence is useless. In summary, this study yields one clean null and one uninterpretable one, not two independent confirmations. A predictor with native uncertainty --- quantile or Bayesian regression, or a conformal wrapper --- would remove the mismatch and is the concrete design a follow-up should adopt.

% PROVENANCE: Figure 3 (eval_risk_coverage.png); Table 5. Source: kind=eval_selective, kind=eval_coverage.
\subsection{Robustness}

Conclusions do not depend on the filtering choice. Repeating the LOUO analysis on unfiltered features and at an 8 Hz cutoff leaves knot-tying positive throughout (+1.18 to +1.69 GRS points) and suturing and needle-passing non-significant throughout; knot-tying's interval excludes zero in four of the six configuration-by-filter combinations, failing at 8 Hz (lower bound -0.26). Widening the ridge penalty grid in a second configuration leaves knot-tying unchanged and moves suturing and needle-passing toward the reference, because the inner search then selects maximal shrinkage; both configurations are reported.

The zero-shot model is close to, but not exactly, deterministic: three repetitions of an identical request agree at Spearman 0.999 with a mean absolute score difference of 0.022-0.024 on a 0-4 scale, and only 19-21\% of trials return a bit-identical score (\texttt{addendum\_jev\_repeatability}). Averaging the three runs changes every headline figure by at most 0.01 GRS points.

% PROVENANCE: Source: kind=headline_mae_vs_reference and kind=jev_headline_mae_vs_reference across
\texttt{feature\_set} values \texttt{smoothed\_5hz}, \texttt{raw} and \texttt{smoothed\_8hz}; \texttt across both configuration files.*

---

\section{Discussion}

The skill label here is a property of the surgeon, constant across their trials. Under the permissive protocol every test surgeon also appears in training, so a model that recognises an individual can retrieve that individual's label without representing skill, and will be scored correct for it. Our results are consistent with this: the protocol alone accounts for at least a third of reported accuracy, and the features identify the operator far above chance. We have not shown the classifier internally relies on identity, and a reader may hold the weaker view that both results reflect some third property of individuals; our data do not separate those accounts. Either way, a figure obtained with the surgeon present in training does not estimate performance on a new surgeon, the only situation in which such a tool would be deployed. Any dataset predicting a person-level label from repeated measurements of the same people admits the same shortcut, though we measured it only here.

What survives is modest: one task of three beats a reference using no motion information, by roughly a quarter to a third of the per-task standard deviation, and there a single-feature regression on duration matches both models with heavily overlapping intervals. The zero-shot model's output correlates with duration at -0.81 to -0.91, so the convergence is not an artefact of fitting. We disclose that its rubric, written before any trial was scored but after we had seen which features carried signal, refers to completion time and pauses; this is a plausible route by which it came to track duration and we cannot exclude it.

For deployment the relevant finding is that confidence does not track correctness on unfamiliar operators. Both paths are overconfident and score worse than a constant forecast; a system confidently wrong on new surgeons is worse for triage than one uncertain everywhere, because the errors it passes through unreviewed are the ones it will not flag. Neither path's confidence identified deferrable trials, but the two nulls are not equal evidence: one is clean, one uninterpretable by construction, and both under-powered. A predictor with native uncertainty, through quantile or conformal methods, would remove the pairing defect and is the improvement a follow-up should adopt.

Three cheap changes would have caught this: report the held-out-surgeon protocol as primary with the paired difference beside it; anchor claims to an uninformative reference, since a positive correlation can accompany larger error than predicting the training mean; and publish a subject-identifiability diagnostic alongside a one-feature baseline. We note the tension in recommending protocol discipline from a study that substituted its own headline metric and added its central diagnostic post hoc; we have labelled both rather than presenting the final analysis as the planned one.

This is a single-dataset case study. Eight surgeons, seven on one task, make every interval crude, and those intervals express which surgeons were evaluated rather than refitting variability or the sampling of these eight. We examined one baseline and one commercial model at one pinned version with one rubric. We reproduced no published system, so our claims concern the protocol's structure and our own demonstration of it, not the validity of specific prior results. The confound and the calibration failure are demonstrated here, not established as general, and both should be checked on a second, larger dataset before being treated as properties of the field.


%% Figures and tables are collected after the text; move them inline if the class prefers.
\begin{figure*}[htbp]
\centering
\includegraphics[width=1.0\textwidth]{fig1_protocol_gap.png}
\caption{The protocol gap, its explanation, and what survives it. \textbf{A}: three-way skill classification accuracy under each protocol, with the majority-class rate marked. \textbf{B}: accuracy of identifying which surgeon performed a held-out trial, from the same features, with chance marked. \textbf{C}: reduction in GRS error relative to a feature-free reference under leave-one-user-out. All intervals are 95\%, resampling whole surgeons.}
\label{fig:protocol_gap}
\end{figure*}

\begin{figure*}[htbp]
\centering
\includegraphics[width=1.0\textwidth]{eval_reliability.png}
\caption{Reliability diagrams. Each path is judged against the target it was given, so the panels are not a head-to-head comparison. Points below the diagonal are overconfident.}
\label{fig:reliability}
\end{figure*}

\begin{figure*}[htbp]
\centering
\includegraphics[width=0.92\textwidth]{eval_risk_coverage.png}
\caption{Risk--coverage curves. The coloured curve should fall below the dotted random-order reference as coverage drops if confidence is informative; the dashed line is the best any ordering could achieve.}
\label{fig:risk_coverage}
\end{figure*}

\begin{figure*}[htbp]
\centering
\includegraphics[width=1.0\textwidth]{eval_bland_altman_LOUO.png}
\caption{Bland--Altman agreement on held-out surgeons. The downward slope is partly intrinsic for any weak predictor and is not independent evidence.}
\label{fig:bland_altman}
\end{figure*}


\begin{table}[htbp]
\centering
\small
\resizebox{\textwidth}{!}{%
\begin{tabular}{lllllllll}
\toprule
\textbf{task} & \textbf{LOSO accuracy} & \textbf{LOUO accuracy} & \textbf{LOSO - LOUO} & \textbf{majority class} & \textbf{LOSO kappa} & \textbf{LOUO kappa} & \textbf{surgeon ID} & \textbf{chance} \\
\midrule
Suturing & 0.92 [0.85, 0.97] & 0.44 [0.25, 0.63] & 0.49 [0.35, 0.62] & 0.487 & +0.77 & -0.19 & 0.87 [0.77, 0.97] & 0.125 \\
Knot-Tying & 0.97 [0.92, 1.00] & 0.31 [0.08, 0.58] & 0.67 [0.41, 0.89] & 0.444 & +0.98 & +0.13 & 0.83 [0.67, 0.97] & 0.125 \\
Needle-Passing & 0.93 [0.85, 1.00] & 0.39 [0.16, 0.69] & 0.54 [0.27, 0.71] & 0.393 & +0.80 & -0.46 & 0.93 [0.83, 1.00] & 0.143 \\
\bottomrule
\end{tabular}
}
\caption{The protocol gap, and surgeon identifiability. Three-way skill classification (novice / intermediate / expert). Brackets are 95\% intervals from resampling whole surgeons. The LOSO - LOUO column is bootstrapped paired on the same resampled surgeons. 'Surgeon ID' replaces the target with surgeon identity, leaving features, model and folds unchanged; it is a post-hoc diagnostic. Every LOUO interval spans its task's majority-class rate.}
\label{tab:table2}
\end{table}

\begin{table}[htbp]
\centering
\small
\resizebox{\textwidth}{!}{%
\begin{tabular}{lllllll}
\toprule
\textbf{scheme} & \textbf{task} & \textbf{reference MAE} & \textbf{method} & \textbf{model MAE} & \textbf{reduction [95\%]} & \textbf{beats reference} \\
\midrule
LOUO & Suturing & 4.71 & 16 features + ridge & 4.56 & 0.15 [-1.04, 1.33] & no \\
 &  &  & Jev, zero-shot & 4.77 & -0.05 [-1.05, 0.96] & no \\
 &  &  & task time alone & 4.61 & 0.10 [-0.72, 0.83] & no \\
LOUO & Knot-Tying & 5.10 & 16 features + ridge & 3.56 & 1.54 [0.65, 2.52] & \textbf{yes} \\
 &  &  & Jev, zero-shot & 2.80 & 2.30 [1.14, 3.45] & \textbf{yes} \\
 &  &  & task time alone & 3.22 & 1.88 [0.94, 2.89] & \textbf{yes} \\
LOUO & Needle-Passing & 4.41 & 16 features + ridge & 4.91 & -0.50 [-1.21, -0.04] & no \\
 &  &  & Jev, zero-shot & 4.65 & -0.24 [-0.79, 0.31] & no \\
 &  &  & task time alone & 4.70 & -0.29 [-0.88, 0.11] & no \\
LOSO & Suturing & 4.35 & 16 features + ridge & 3.20 & 1.15 [0.47, 1.92] & \textbf{yes} \\
 &  &  & Jev, zero-shot & 4.02 & 0.33 [-0.70, 1.51] & no \\
 &  &  & task time alone & 4.01 & 0.34 [-0.32, 1.01] & no \\
LOSO & Knot-Tying & 4.66 & 16 features + ridge & 2.88 & 1.78 [1.00, 2.58] & \textbf{yes} \\
 &  &  & Jev, zero-shot & 2.59 & 2.08 [1.07, 3.09] & \textbf{yes} \\
 &  &  & task time alone & 2.84 & 1.82 [1.00, 2.68] & \textbf{yes} \\
LOSO & Needle-Passing & 4.22 & 16 features + ridge & 3.79 & 0.43 [0.11, 0.76] & \textbf{yes} \\
 &  &  & Jev, zero-shot & 4.27 & -0.05 [-0.51, 0.53] & no \\
 &  &  & task time alone & 4.40 & -0.18 [-0.48, 0.15] & no \\
\bottomrule
\end{tabular}
}
\caption{GRS error against the feature-free reference. The reference predicts each held-out trial with the mean GRS of its own training fold. A positive reduction means lower error than that reference; 'beats reference' requires the lower bound to exceed zero, a rule fixed before any interval was computed. Jev's ordinal score is placed on the GRS scale by a line fitted within each training fold.}
\label{tab:table3}
\end{table}

\begin{table}[htbp]
\centering
\small
\resizebox{\textwidth}{!}{%
\begin{tabular}{lllllll}
\toprule
\textbf{scheme} & \textbf{task} & \textbf{reference MAE} & \textbf{method} & \textbf{model MAE} & \textbf{reduction [95\%]} & \textbf{beats reference} \\
\midrule
LOUO & Suturing & 4.71 & 16 features + ridge & 4.56 & 0.15 [-1.04, 1.33] & no \\
 &  &  & Jev, zero-shot & 4.77 & -0.05 [-1.05, 0.96] & no \\
 &  &  & task time alone & 4.61 & 0.10 [-0.72, 0.83] & no \\
LOUO & Knot-Tying & 5.10 & 16 features + ridge & 3.56 & 1.54 [0.65, 2.52] & \textbf{yes} \\
 &  &  & Jev, zero-shot & 2.80 & 2.30 [1.14, 3.45] & \textbf{yes} \\
 &  &  & task time alone & 3.22 & 1.88 [0.94, 2.89] & \textbf{yes} \\
LOUO & Needle-Passing & 4.41 & 16 features + ridge & 4.91 & -0.50 [-1.21, -0.04] & no \\
 &  &  & Jev, zero-shot & 4.65 & -0.24 [-0.79, 0.31] & no \\
 &  &  & task time alone & 4.70 & -0.29 [-0.88, 0.11] & no \\
LOSO & Suturing & 4.35 & 16 features + ridge & 3.20 & 1.15 [0.47, 1.92] & \textbf{yes} \\
 &  &  & Jev, zero-shot & 4.02 & 0.33 [-0.70, 1.51] & no \\
 &  &  & task time alone & 4.01 & 0.34 [-0.32, 1.01] & no \\
LOSO & Knot-Tying & 4.66 & 16 features + ridge & 2.88 & 1.78 [1.00, 2.58] & \textbf{yes} \\
 &  &  & Jev, zero-shot & 2.59 & 2.08 [1.07, 3.09] & \textbf{yes} \\
 &  &  & task time alone & 2.84 & 1.82 [1.00, 2.68] & \textbf{yes} \\
LOSO & Needle-Passing & 4.22 & 16 features + ridge & 3.79 & 0.43 [0.11, 0.76] & \textbf{yes} \\
 &  &  & Jev, zero-shot & 4.27 & -0.05 [-0.51, 0.53] & no \\
 &  &  & task time alone & 4.40 & -0.18 [-0.48, 0.15] & no \\
\bottomrule
\end{tabular}
}
\caption{GRS error against the feature-free reference. The reference predicts each held-out trial with the mean GRS of its own training fold. A positive reduction means lower error than that reference; 'beats reference' requires the lower bound to exceed zero, a rule fixed before any interval was computed. Jev's ordinal score is placed on the GRS scale by a line fitted within each training fold.}
\label{tab:table3}
\end{table}

\begin{table}[htbp]
\centering
\small
\resizebox{\textwidth}{!}{%
\begin{tabular}{llllllllll}
\toprule
\textbf{path} & \textbf{scheme} & \textbf{judged against} & \textbf{mean confidence} & \textbf{accuracy} & \textbf{ECE [95\%]} & \textbf{MCE} & \textbf{Brier} & \textbf{constant Brier} & \textbf{Brier skill [95\%]} \\
\midrule
baseline & LOUO & self-declared skill (N/I/E) & 0.755 & 0.379 & 0.377 [0.237, 0.541] & 0.53 & 0.875 & 0.647 & -0.351 [-0.688, -0.092] \\
baseline & LOSO & self-declared skill (N/I/E) & 0.837 & 0.942 & 0.105 [0.057, 0.154] & 0.37 & 0.143 & 0.647 & 0.778 [0.698, 0.856] \\
jev & both & level implied by GRS & 0.696 & 0.301 & 0.395 [0.317, 0.494] & 0.55 & 1.021 & 0.713 & -0.433 [-0.588, -0.242] \\
\bottomrule
\end{tabular}
}
\caption{Calibration. Each path is judged against the target it was given, so the rows are not a head-to-head comparison. The constant forecast assigns every trial the observed class frequencies; a negative Brier skill score means the path is worse than that constant. Five equal-width bins, 103 pooled trials. Jev is never trained, so it has a single row. Per-bin counts are logged as \texttt{addendum\_reliability\_bin}; mean confidence is taken from \texttt{addendum\_mean\_confidence}, which postdates the calibration run.}
\label{tab:table4}
\end{table}

\begin{table}[htbp]
\centering
\small
\resizebox{\textwidth}{!}{%
\begin{tabular}{llllllll}
\toprule
\textbf{path} & \textbf{scheme} & \textbf{overall MAE} & \textbf{AURC} & \textbf{oracle AURC} & \textbf{E-AURC} & \textbf{usefulness [95\%]} & \textbf{informative} \\
\midrule
baseline & LOUO & 4.31 & 4.17 & 1.86 & 2.32 & 0.13 [-0.63, 0.92] (5.5\% of the gap to ideal) & no \\
jev & LOUO & 4.05 & 3.89 & 1.96 & 1.93 & 0.16 [-0.42, 0.61] (7.4\% of the gap to ideal) & no \\
jev (max prob) & LOUO & 4.05 & 3.90 & 1.96 & 1.94 & 0.15 [-0.45, 0.63] (7.2\% of the gap to ideal) & no \\
baseline & LOSO & 3.25 & 3.20 & 1.42 & 1.78 & 0.05 [-0.40, 0.67] (2.5\% of the gap to ideal) & no \\
jev & LOSO & 3.59 & 3.51 & 1.64 & 1.87 & 0.08 [-0.40, 0.41] (3.9\% of the gap to ideal) & no \\
jev (max prob) & LOSO & 3.59 & 3.53 & 1.64 & 1.89 & 0.06 [-0.44, 0.40] (3.1\% of the gap to ideal) & no \\
path & scheme & 100\% & 80\% & 60\% & 40\% \\
--- & --- & --- & --- & --- & --- \\
baseline & LOUO & 4.31 (+0.46) & 4.50 (+0.23) & 4.37 (+0.34) & 4.15 (+0.39) \\
jev & LOUO & 4.05 (+0.72) & 4.27 (+0.55) & 4.22 (+0.71) & 3.76 (+1.11) \\
baseline & LOSO & 3.25 (+1.17) & 3.07 (+1.36) & 3.28 (+1.32) & 3.09 (+1.51) \\
jev & LOSO & 3.59 (+0.84) & 3.84 (+0.60) & 3.82 (+0.70) & 3.40 (+1.16) \\
\bottomrule
\end{tabular}
}
\caption{Selective prediction. Usefulness is the overall MAE minus the area under the risk-coverage curve: positive means ordering by confidence lowers the error of the trials kept, relative to deferring at random. 'Informative' requires the lower bound to exceed zero. Mean absolute GRS error of the trials retained at each coverage, with the reduction against the feature-free reference on those same trials in brackets. Coverage levels were fixed in advance.}
\label{tab:table5}
\end{table}


\bibliographystyle{plainnat}
\bibliography{references}

\end{document}
